\documentclass[10pt,twocolumn,letterpaper]{article}

% ── Paquetes ─────────────────────────────────────────────────────────────────
\usepackage[utf8]{inputenc}
\usepackage[T1]{fontenc}
\usepackage[spanish,es-nodecimaldot]{babel}
\usepackage{times}
\usepackage{geometry}
\usepackage{graphicx}
\usepackage{amsmath}
\usepackage{booktabs}
\usepackage{tabularx}
\usepackage{array}
\usepackage{float}
\usepackage{caption}
\usepackage{cite}
\usepackage{url}
\usepackage{microtype}
\usepackage{parskip}
\usepackage{titlesec}
\usepackage{fancyhdr}
\usepackage[hidelinks]{hyperref}
\usepackage{enumitem}
\renewcommand{\baselinestretch}{1.05}

% ── Geometría de página ───────────────────────────────────────────────────────
\geometry{
  top=2.0cm, bottom=2.2cm,
  left=1.8cm, right=1.8cm,
  columnsep=0.65cm
}

% ── Encabezado/pie ─────────────────────────────────────────────────────────────
\pagestyle{fancy}
\fancyhf{}
\renewcommand{\headrulewidth}{0pt}
\fancyfoot[C]{\thepage}

% ── Formato de secciones ──────────────────────────────────────────────────────
\titleformat{\section}[block]
  {\normalfont\normalsize\bfseries\centering}
  {\Roman{section}.}
  {0.5em}
  {\MakeUppercase}
\titlespacing*{\section}{0pt}{9pt}{4pt}

\titleformat{\subsection}[runin]
  {\normalfont\normalsize\itshape}
  {\Alph{subsection}.}
  {0.5em}
  {}[.]
\titlespacing*{\subsection}{0pt}{5pt}{4pt}

% ── Captions ──────────────────────────────────────────────────────────────────
\captionsetup{font=small, labelfont=bf, labelsep=period, justification=centering}
\captionsetup[table]{position=top, aboveskip=3pt, belowskip=2pt}
\captionsetup[figure]{position=bottom, aboveskip=2pt, belowskip=0pt}

% ── Espaciado ─────────────────────────────────────────────────────────────────
\setlength{\parindent}{0pt}
\setlength{\parskip}{4.5pt plus 1pt}

% ─────────────────────────────────────────────────────────────────────────────
\begin{document}

% ── Encabezado institucional ──────────────────────────────────────────────────
\twocolumn[{%
\begin{center}
  {\small Instituto Politécnico Nacional \quad|\quad Escuela Superior de Cómputo \quad TT\,2026--A062}\\[4pt]
  {\LARGE\bfseries FINARA: Prototipo de Aplicación para Finanzas\\[2pt]
   Personales con Escaneo de Tickets usando\\[2pt]
   Reconocimiento Óptico de Caracteres}\\[6pt]
  {\normalsize
   Acevedo Martínez Armando,\quad Cadenas Acevedo Jesús Alejandro\\
   Directores: Prof.\ Ulises Vélez Saldaña,\quad Dra.\ María de Lourdes Salinas Callejas}\\[2pt]
  {\small Escuela Superior de Cómputo, Instituto Politécnico Nacional, México D.F.}\\[8pt]
\end{center}

\noindent\rule{\textwidth}{0.4pt}\\[4pt]
\noindent\textit{\textbf{Resumen}~---
\textbf{
Se presenta FINARA, un sistema integral de gestión de finanzas personales compuesto por una aplicación
móvil (Flutter), una plataforma web (Next.js/React) y un servidor REST (Node.js/Express) con
persistencia en MySQL. El sistema automatiza el registro de gastos a partir de tickets físicos y
digitales mediante un \emph{pipeline} híbrido de OCR que combina Tesseract y Google Cloud Vision, con
categorización automática por palabras clave. Un módulo de recomendaciones financieras personalizadas
emplea Claude Haiku como motor de análisis con IA. La evaluación sobre 30 tickets físicos alcanzó una
tasa de detección del campo total del 86.3\,\%. En una batería de 33 casos de prueba se obtuvo el
100\,\% de aprobación en seguridad y compatibilidad. La encuesta de usabilidad aplicada a 20
participantes arrojó un NPS de~60 con promedio de 8.9/10.
}}\\[4pt]
\noindent\textit{\textbf{Palabras~clave}~--- Finanzas personales, reconocimiento óptico de caracteres,
aplicación móvil, arquitectura cliente-servidor, Flutter, Node.js, MySQL, Google Cloud Vision.}\\[4pt]
\noindent\rule{\textwidth}{0.4pt}\\[6pt]
}]


% ══════════════════════════════════════════════════════════════════════════════
\section{Introducción}
% ══════════════════════════════════════════════════════════════════════════════

Según la Encuesta Nacional sobre Salud Financiera (ENSAFI) 2023 del INEGI, solo el 53.2\,\% de los
mexicanos lleva algún tipo de registro de sus gastos y únicamente el 32.8\,\% cumple con dicho
registro de manera continua \cite{inegi2024}. El mismo estudio revela que el 45.9\,\% de la
población reporta que casi nunca o nunca le sobra dinero al final del mes, evidenciando capacidad
de ahorro limitada. Adicionalmente, un alto porcentaje de personas en múltiples países se muestra
preocupado por su situación financiera futura \cite{statista2025}.

A pesar de que las instalaciones de aplicaciones de finanzas personales aumentaron un 50\,\% en
2023 respecto al año anterior, la retención de usuarios cae drásticamente a partir del séptimo día
de uso \cite{adjust2023}. Esto indica que las herramientas disponibles no logran sostener el cambio
de hábitos. Estudios de la OCDE/INFE (2023) señalan que una mayor alfabetización financiera
apoyada por herramientas digitales está asociada con comportamientos más responsables y mayor
capacidad de ahorro \cite{OECD_INFE2023}.

Esta brecha entre necesidad y adopción tecnológica motivó el desarrollo de FINARA: una plataforma
que reduce la carga manual del registro financiero mediante automatización basada en Reconocimiento
Óptico de Caracteres (OCR), categorización inteligente de gastos y recomendaciones financieras
generadas por inteligencia artificial. La solución está orientada a jóvenes y adultos sin
conocimientos contables avanzados, apoyándose en la disponibilidad del 81.4\,\% de la población
mexicana a dispositivos móviles \cite{EncINEGI}.

El objetivo general fue \textbf{desarrollar una aplicación móvil y plataforma web interactiva para
la gestión de finanzas personales que registre automáticamente los gastos de tickets físicos y
digitales, categorizándolos para facilitar el análisis y seguimiento del presupuesto del usuario}.
Los objetivos específicos incluyeron: módulo de escaneo de tickets físicos y digitales, módulo de
categorización automática de gastos, módulo de alertas y notificaciones, generación de informes
exportables y visualización de análisis con gráficas interactivas.


% ══════════════════════════════════════════════════════════════════════════════
\section{Estado del Arte}
% ══════════════════════════════════════════════════════════════════════════════

Se analizaron las principales aplicaciones existentes de gestión financiera con funcionalidades
similares a las propuestas. La Tabla~\ref{tab:estadoarte} sintetiza las características más
relevantes para la comparativa.

\begin{table}[H]
  \centering
  \small
  \caption{Comparativa de aplicaciones de finanzas personales}
  \label{tab:estadoarte}
  \renewcommand{\arraystretch}{1.1}
  \begin{tabular}{p{1.5cm}p{1.4cm}p{1.1cm}p{1.2cm}p{1.2cm}}
    \toprule
    \textbf{App} & \textbf{OCR} & \textbf{IA} & \textbf{Web+Móvil} & \textbf{Gratis} \\
    \midrule
    Expensify & Sí & Parcial & Sí & No \\
    Veryfi    & Sí & Sí      & Sí & No \\
    Mint      & No & Sí      & Sí & Sí \\
    Shoeboxed & Sí & No      & Sí & No \\
    Fyle      & Sí & Sí      & Sí & No \\
    \textbf{FINARA} & \textbf{Sí} & \textbf{Sí} & \textbf{Sí} & \textbf{Sí} \\
    \bottomrule
  \end{tabular}
\end{table}

\textbf{Expensify} destaca por su OCR y generación de reportes, pero su interfaz está orientada al
entorno empresarial y requiere validación manual de categorías. \textbf{Veryfi} combina OCR con
aprendizaje automático para categorización en tiempo real, aunque su versión gratuita es limitada
y presenta curva de aprendizaje alta. \textbf{Mint} ofrece integración bancaria y categorización
inteligente, pero carece de módulo de escaneo de tickets. \textbf{Shoeboxed} permite digitalizar
recibos, pero su clasificación de gastos no es precisa en español. \textbf{Fyle} logra alta
precisión con automatización completa, pero su modelo de negocio corporativo restringe su adopción
por usuarios individuales.

FINARA se diferencia por ser la única solución completamente gratuita y de código abierto que
integra OCR híbrido con respaldo en la nube, recomendaciones por IA y soporte completo para
tickets en español, orientada específicamente a la problemática financiera del usuario mexicano.


% ══════════════════════════════════════════════════════════════════════════════
\section{Metodología}
% ══════════════════════════════════════════════════════════════════════════════

\subsection{Arquitectura del sistema}

FINARA adopta una arquitectura cliente-servidor de tres capas \cite{ArquitecturaCapas} que separa
la interfaz de usuario, la lógica de negocio y la gestión de datos. Los tres componentes se
comunican mediante una API REST autenticada con JSON Web Tokens (JWT), garantizando el aislamiento
entre cuentas de usuario y la coherencia de datos entre plataformas (Fig.~\ref{fig:arq}).

\begin{figure}[H]
  \centering
  \includegraphics[width=\columnwidth]{images/ArquitecturaCapasFinara.JPG}
  \caption{Arquitectura de tres capas del sistema FINARA.}
  \label{fig:arq}
\end{figure}

Los tres componentes son:

\begin{itemize}[noitemsep,topsep=2pt]
  \item \textbf{Aplicación móvil (Flutter):} 9 módulos funcionales orientados al registro
    ágil de gastos, escaneo de tickets físicos mediante la cámara del dispositivo, gestión
    de presupuestos, deudas y alertas automáticas. Compatible con Android 10 o superior.
  \item \textbf{Plataforma web (Next.js/React):} 6 módulos enfocados en análisis financiero
    detallado, visualización de gráficas comparativas, generación de reportes exportables en
    PDF y Excel, y recomendaciones personalizadas por período.
  \item \textbf{Backend REST (Node.js/Express):} 11 módulos con más de 60 endpoints
    documentados, persistencia en MySQL (12 tablas normalizadas a 3FN), tareas programadas
    (cron~jobs) para recordatorios diarios y rollover mensual automático de presupuestos e
    ingresos recurrentes.
\end{itemize}

La seguridad se implementó mediante cifrado de contraseñas con bcrypt (factor~12) y validación JWT
en todos los endpoints protegidos. La base de datos MySQL fue diseñada con relaciones de integridad
referencial y eliminación en cascada.

La metodología de desarrollo adoptada fue una combinación de Extreme Programming (XP) y Kanban
\cite{Metodologias,ManifestoAgil}, con iteraciones cortas de dos semanas, integración continua y
revisión constante de requerimientos con los directores del trabajo.

\subsection{Pipeline de OCR híbrido}

El módulo de reconocimiento óptico de caracteres es el componente técnicamente más elaborado del
sistema. Se diseñó un \emph{pipeline} de seis etapas (Fig.~\ref{fig:ocr}) que combina
procesamiento local con servicios externos, optimizando el costo de operación:

\begin{enumerate}[noitemsep,topsep=2pt]
  \item \textbf{Preprocesamiento:} Corrección de iluminación, reescalado, corrección de
    perspectiva (\textit{deskew}) y binarización adaptativa con OpenCV.
  \item \textbf{Extracción con Tesseract:} OCR con modelo LSTM en modo de bloque de texto
    (PSM~6), con soporte para español e inglés \cite{Tesseract}.
  \item \textbf{Evaluación de confianza:} Métrica ponderada que combina la confianza reportada
    por Tesseract con la presencia detectada de campos clave: total, productos, tienda y fecha.
  \item \textbf{Compuertas de activación:} Si la confianza es inferior a 0.72 o el número de
    palabras reconocidas es menor a 15, el flujo escala automáticamente al servicio externo.
  \item \textbf{Respaldo externo:} Google Cloud Vision extrae el texto completo del documento;
    Claude Haiku lo estructura en JSON con los campos requeridos \cite{GoogleCloudVisionAPI}.
    A diferencia de alternativas como AWS Textract o Azure Document Intelligence, Cloud Vision
    está diseñada para texto en imágenes tomadas con cámara en condiciones no controladas, que
    es precisamente el caso de los tickets físicos de FINARA.
  \item \textbf{Categorización automática:} Los productos detectados reciben una categoría
    sugerida mediante un clasificador por palabras clave y similitud con el historial de
    compras previas del usuario.
\end{enumerate}

\begin{figure}[H]
  \centering
  \includegraphics[width=\columnwidth]{images/Flujo-OCR2.png}
  \caption{Diagrama del \textit{pipeline} híbrido de OCR.}
  \label{fig:ocr}
\end{figure}

Este diseño permite procesar los tickets de alta calidad localmente, sin costo de API externa, y
escalar únicamente los casos de mayor dificultad a servicios de visión artificial en la nube. Un
sistema de caché con vigencia de 15~días evita llamadas redundantes al modelo de IA para el mismo
período financiero.

\subsection{Módulo de recomendaciones con IA}

Se implementó un módulo de recomendaciones financieras personalizadas que utiliza Claude Haiku como
motor de análisis. Para cada período consultado, el sistema recopila el contexto financiero del
usuario (ingresos, distribución de gastos, estado de presupuestos y deudas) y genera un análisis
estructurado con cuatro secciones: resumen textual, riesgos identificados, recomendaciones
accionables y acciones priorizadas. Ante una falla del servicio externo, el sistema retorna
automáticamente recomendaciones generadas por reglas locales, garantizando la disponibilidad del
módulo en todo momento.

\subsection{Stack tecnológico}

La selección de tecnologías se realizó priorizando herramientas de código abierto, comunidad
activa y compatibilidad multiplataforma. La Tabla~\ref{tab:stack} resume los componentes
principales del sistema.

\begin{table}[H]
  \centering
  \small
  \caption{Stack tecnológico de FINARA}
  \label{tab:stack}
  \renewcommand{\arraystretch}{1.15}
  \begin{tabular}{p{2.2cm}p{4.4cm}}
    \toprule
    \textbf{Componente}   & \textbf{Rol en el sistema} \\
    \midrule
    Flutter / Dart        & App móvil multiplataforma (Android) \\
    Next.js / React       & Plataforma web con SSR y gestión de componentes \\
    Node.js + Express     & Backend REST: lógica de negocio y orquestación \\
    MySQL                 & Persistencia relacional (12 tablas, 3FN) \\
    Python 3 + OpenCV     & Preprocesamiento e invocación de Tesseract \\
    Tesseract 5 (LSTM)    & OCR local sin costo de API \\
    Google Cloud Vision   & OCR en la nube para tickets difíciles \\
    Claude Haiku (Anthropic) & Estructuración JSON y recomendaciones IA \\
    JWT + bcrypt          & Autenticación y cifrado de contraseñas \\
    \bottomrule
  \end{tabular}
\end{table}

Flutter fue seleccionado frente a React Native por su rendimiento nativo, widget unificado y
compilación AOT, que garantiza una experiencia fluida en dispositivos de gama baja \cite{WhyFlutter}.
Next.js fue preferido sobre alternativas como Remix o SvelteKit por su ecosistema maduro, su
soporte nativo para renderizado del lado del servidor (SSR) y la facilidad de integración con
bibliotecas de visualización como ECharts \cite{NextJSOfficial}. MySQL se eligió sobre PostgreSQL
y MongoDB por su madurez, rendimiento en consultas analíticas y la familiaridad del equipo con el
modelo relacional \cite{WhyMySQL}.

\subsection{Modelo de categorización de gastos}

La asignación de categorías a los productos extraídos del ticket se basa en el
\textbf{modelo de sobres de Dave Ramsey}, adaptado a cuatro categorías principales: Alimentos,
Entretenimiento, Transporte y Otros. El clasificador aplica una búsqueda por palabras clave sobre
el nombre del producto y, si no hay coincidencia directa, consulta el historial de compras previas
del usuario para inferir la categoría por similitud. Los presupuestos se gestionan en forma de
límites mensuales por categoría, con indicadores visuales de tres niveles: seguro (consumo menor
al 75\,\% del límite), advertencia (75\,\%--100\,\%) y peligro/excedido (superación del límite).
El módulo de deudas implementa el cálculo de cuota mensual por amortización francesa, con soporte
para pagos ordinarios y abonos extraordinarios a capital.


% ══════════════════════════════════════════════════════════════════════════════
\section{Resultados}
% ══════════════════════════════════════════════════════════════════════════════

\subsection{Evaluación del módulo OCR}

Se evaluaron \textbf{30 tickets físicos} de establecimientos como Walmart, Soriana y Farmacias del
Ahorro, clasificados en tres niveles de calidad según iluminación, enfoque y estado del
comprobante. La Tabla~\ref{tab:ocr} muestra la distribución de procesamiento y las tasas de
detección de los campos clave.

\begin{table}[H]
  \centering
  \small
  \caption{Métricas de extracción OCR por calidad de imagen (30 tickets)}
  \label{tab:ocr}
  \begin{tabular}{lcccc}
    \toprule
    \textbf{Condición} & \textbf{Tck.} & \textbf{Tess.} & \textbf{GCV} & \textbf{Total det.} \\
    \midrule
    Alta calidad  & 12 &  7 & 5  & 95.8\,\% \\
    Calidad media & 10 &  2 & 8  & 87.0\,\% \\
    Baja calidad  &  8 &  0 & 8  & 73.0\,\% \\
    \midrule
    \textbf{Total} & \textbf{30} & \textbf{9} & \textbf{21} & \textbf{86.3\,\%} \\
    \bottomrule
  \end{tabular}
\end{table}

Google Cloud Vision (GCV) supera a Tesseract en todos los campos clave con una ventaja promedio
de aproximadamente 35 puntos porcentuales (Tabla~\ref{tab:campos}), justificando el diseño
híbrido del \emph{pipeline}.

\begin{table}[H]
  \centering
  \small
  \caption{Comparativa de detección de campos clave: Tesseract vs.\ GCV}
  \label{tab:campos}
  \begin{tabular}{lcc}
    \toprule
    \textbf{Campo clave} & \textbf{Tesseract} & \textbf{Cloud Vision} \\
    \midrule
    Total del ticket & 65.2\,\% & 93.5\,\% \\
    Productos        & 58.8\,\% & 88.7\,\% \\
    Nombre de tienda & 52.4\,\% & 87.3\,\% \\
    Fecha            & 44.4\,\% & 82.9\,\% \\
    \bottomrule
  \end{tabular}
\end{table}

\subsection{Precisión del clasificador automático}

Para evaluar el clasificador se probaron \textbf{150 productos} distribuidos en cuatro categorías.
La exactitud global fue del \textbf{62.0\,\%} y la puntuación F1 macro promedio fue de
\textbf{57.7\,\%}. La Tabla~\ref{tab:metricas} detalla las métricas por categoría.

\begin{table}[H]
  \centering
  \small
  \caption{Métricas del clasificador de categorías (150 productos)}
  \label{tab:metricas}
  \begin{tabular}{lccc}
    \toprule
    \textbf{Categoría} & \textbf{Precisión} & \textbf{Exhaustiv.} & \textbf{F1} \\
    \midrule
    Alimentos       & 65.8\,\% & 83.3\,\% & 73.5\,\% \\
    Entretenimiento & 53.3\,\% & 40.0\,\% & 45.7\,\% \\
    Transporte      & 66.7\,\% & 56.0\,\% & 60.9\,\% \\
    Otros           & 55.3\,\% & 46.7\,\% & 50.6\,\% \\
    \midrule
    \textbf{Macro prom.} & \textbf{60.3\,\%} & \textbf{56.5\,\%} & \textbf{57.7\,\%} \\
    \bottomrule
  \end{tabular}
\end{table}

La categoría \textit{Alimentos} concentra la mayoría de las predicciones con alta exhaustividad
(83.3\,\%) pero baja precisión (65.8\,\%), absorbiendo productos de otras categorías. \textit{Entretenimiento}
es la más afectada con F1 de apenas 45.7\,\%, principalmente por la asignación incorrecta de
suscripciones digitales y accesorios de videojuegos. Estos resultados establecen la línea base
para mejorar el clasificador en iteraciones futuras.

\subsection{Validación mediante casos de prueba}

Se diseñaron y ejecutaron \textbf{33 casos de prueba de caja negra} en \textbf{9 dispositivos
físicos Android} de distintas gamas (alta, media-alta, media y baja), agrupados en seis
categorías (Tabla~\ref{tab:pruebas}).

\begin{table}[H]
  \centering
  \small
  \caption{Resultados de pruebas por categoría (33 casos)}
  \label{tab:pruebas}
  \begin{tabular}{lcccr}
    \toprule
    \textbf{Categoría} & \textbf{Total} & \textbf{Pasó} & \textbf{Falló} & \textbf{\% Aprob.} \\
    \midrule
    Funcionalidad  & 21 & 13 &  8 & 61.9\,\% \\
    Seguridad      &  3 &  3 &  0 & 100.0\,\% \\
    Usabilidad     &  4 &  3 &  1 & 75.0\,\% \\
    Portabilidad   &  3 &  2 &  1 & 66.7\,\% \\
    Compatibilidad &  1 &  1 &  0 & 100.0\,\% \\
    Mantenibilidad &  1 &  0 &  1 &   0.0\,\% \\
    \midrule
    \textbf{Total} & \textbf{33} & \textbf{22} & \textbf{11} & \textbf{66.7\,\%} \\
    \bottomrule
  \end{tabular}
\end{table}

Los módulos de seguridad y compatibilidad alcanzaron el 100\,\% de aprobación. La autenticación
sin estado con JWT, el cifrado con bcrypt (factor~12) y la resistencia demostrada a inyecciones
SQL conforman una base de seguridad adecuada para datos financieros sensibles. La compatibilidad
fue verificada con Android 10--14 (API 29--34) en pantallas de 5 a 6.7 pulgadas en resoluciones
de 720p a 1440p. Los 11 casos fallidos generaron un catálogo de 15~defectos; los más
significativos fueron: ausencia de filtros avanzados en el historial de transacciones, pantalla
de estadísticas por categoría no implementada, ciclo de vida incompleto de deudas liquidadas e
inconsistencias en la edición de categorías. De los 15 defectos, 3 fueron corregidos en la misma
iteración, 3 atendidos parcialmente y 9 quedaron programados para la siguiente fase del proyecto.

\subsection{Evaluación con usuarios}

Se aplicó una encuesta de usabilidad a \textbf{20 participantes} en dos etapas: la primera el
30 de abril de 2026 (6 participantes) y la segunda entre el 26 de mayo y el 2 de junio de 2026
(14 participantes adicionales). La Tabla~\ref{tab:encuesta} resume los resultados más relevantes.

\begin{table}[H]
  \centering
  \small
  \caption{Resultados de la encuesta de usabilidad (20 participantes)}
  \label{tab:encuesta}
  \renewcommand{\arraystretch}{1.15}
  \begin{tabular}{p{2.6cm}p{3.0cm}c}
    \toprule
    \textbf{Dimensión} & \textbf{Resultado más frecuente} & \textbf{Frec.} \\
    \midrule
    Diseño visual        & Me encantó, muy atractivo         & 80.0\,\% \\
    Alertas automáticas  & Sí, muy útil                      & 80.0\,\% \\
    Presupuestos         & Muy útil, la usaría               & 60.0\,\% \\
    Registro manual      & Rápido, 15--30 segundos           & 60.0\,\% \\
    Sección de deudas    & Clara, solo dudas menores         & 60.0\,\% \\
    Función OCR          & Funcionó bien / perfecto          & 80.0\,\% \\
    Navegación           & Fácil, se perdió un par de veces  & 65.0\,\% \\
    \midrule
    \multicolumn{2}{l}{\textbf{NPS (Índice de recomendación)}} & \textbf{60} \\
    \multicolumn{2}{l}{\textbf{Promedio de recomendación}} & \textbf{8.9\,/\,10} \\
    \bottomrule
  \end{tabular}
\end{table}

Los puntajes individuales de recomendación de los 20 participantes arrojaron 13 promotores
(puntuación 9--10, 65\,\%), 6 pasivos (puntuación 7--8, 30\,\%) y 1 detractor (5\,\%), dando
un \textbf{NPS de 60}. Este valor se considera ``bueno'' en la industria de aplicaciones móviles
\cite{BBVA_Finanzas}. La funcionalidad más solicitada para versiones futuras fue metas de ahorro
(40\,\%), seguida de exportación en PDF/Excel (30\,\%) y modo oscuro (15\,\%).

Un hallazgo relevante fue la mejora significativa en la percepción del OCR entre etapas: en la
primera ronda (abril), el 66.7\,\% de los participantes reportó que el módulo reconoció muy poco;
en la segunda ronda (mayo--junio), el 80\,\% reportó que funcionó bien o perfectamente. Este
contraste valida las correcciones aplicadas al \emph{pipeline} y posiciona la mejora continua
del clasificador como prioridad.


% ══════════════════════════════════════════════════════════════════════════════
\section{Conclusiones}
% ══════════════════════════════════════════════════════════════════════════════

El sistema FINARA demostró ser técnicamente viable y funcionalmente competitivo en sus módulos
centrales. Los resultados cuantitativos confirman que la base del sistema es sólida: 66.7\,\%
de tasa de aprobación global en casos de prueba, 100\,\% de éxito en seguridad y compatibilidad,
y un NPS de~60 (promedio 8.9/10) en la evaluación con 20 usuarios.

El \emph{pipeline} híbrido de OCR opera correctamente en la extracción de texto con tasas de
detección del 86.3\,\% en el campo total sobre 30 tickets físicos. Sin embargo, la fase de
\textbf{categorización automática de productos} no alcanzó el nivel de precisión requerido
(F1 macro 57.7\,\%), siendo la principal limitación identificada y la prioridad de la siguiente
iteración del proyecto. Específicamente, el clasificador presenta sesgo hacia la categoría
\textit{Alimentos}, requiriendo la incorporación de reglas léxicas adicionales y datos de
entrenamiento más balanceados.

La percepción de los usuarios fue positiva en todas las dimensiones evaluadas. El diseño visual
fue valorado favorablemente por el 100\,\% de los participantes, y el módulo de escaneo
automático de tickets fue identificado como la funcionalidad diferenciadora del sistema. Los
13 defectos restantes en el catálogo constituyen la agenda de trabajo para la siguiente fase,
con prioridad en la categorización OCR, los filtros del historial y el ciclo de vida de las
deudas.

Con este proyecto se demuestra que la integración de OCR híbrido, arquitectura cliente-servidor
de tres capas y recomendaciones generadas por inteligencia artificial constituye una solución
efectiva y accesible para fomentar hábitos financieros saludables en la población mexicana sin
conocimientos contables avanzados.

\subsection{Trabajo futuro}

Las líneas de trabajo identificadas para la siguiente iteración son:

\begin{enumerate}[noitemsep,topsep=2pt]
  \item \textbf{Mejora del clasificador de categorías:} Incorporar reglas léxicas adicionales,
    ampliar el diccionario de palabras clave por categoría y explorar técnicas de aprendizaje
    automático supervisado para superar la línea base de F1 macro 57.7\,\%.
  \item \textbf{Pantalla de estadísticas por categoría:} Implementar la vista de análisis
    comparativo de gastos por categoría con filtros temporales y desglose por establecimiento.
  \item \textbf{Ciclo de vida completo de deudas:} Corregir la lógica del módulo de deudas
    para que al completar el pago total el estado se actualice automáticamente a ``liquidada''.
  \item \textbf{Filtros avanzados en el historial:} Agregar búsqueda por texto, filtro
    combinado por categoría y tipo de movimiento en la vista de historial de transacciones.
  \item \textbf{Metas de ahorro:} Incorporar el módulo de metas financieras, la
    funcionalidad más solicitada por el 40\,\% de los participantes de la encuesta.
  \item \textbf{Refactorización a arquitectura feature-first:} Migrar la gestión de estado
    de la app móvil de \texttt{StatefulWidget}/\texttt{setState} a un patrón reactivo
    (BLoC o Riverpod) para reducir el acoplamiento entre pantallas.
\end{enumerate}


% ══════════════════════════════════════════════════════════════════════════════
\section*{Reconocimientos}
% ══════════════════════════════════════════════════════════════════════════════

Los autores agradecen al Prof. Ulises Vélez Saldaña y a la Dra. María de Lourdes Salinas Callejas
por su guía y dedicación durante el desarrollo del proyecto, a los sinodales por sus observaciones
que enriquecieron el trabajo, y a la Escuela Superior de Cómputo del Instituto Politécnico Nacional
por las facilidades otorgadas. Asimismo, se agradece a los 20 participantes que colaboraron en las
sesiones de evaluación de usabilidad del sistema FINARA.


% ══════════════════════════════════════════════════════════════════════════════
\begingroup
\small
\renewcommand{\refname}{Referencias}
\begin{thebibliography}{9}

\bibitem{inegi2024}
INEGI, ``Encuesta Nacional sobre Salud Financiera (ENSAFI) 2023,''
Instituto Nacional de Estadística y Geografía, México, 2024. [En línea].
Disponible en: \url{https://www.inegi.org.mx/contenidos/saladeprensa/boletines/2024/ENSAFI/ENSAFI.pdf}

\bibitem{statista2025}
Statista, ``La situación financiera personal en el futuro, ¿fuente de preocupación?,''
Statista, 2025. [En línea].
Disponible en: \url{https://es.statista.com/grafico/27448/porcentaje-de-encuestados-que-estan-preocupados-por-su-futuro-financiero/}

\bibitem{adjust2023}
Adjust, ``El uso de las aplicaciones de finanzas continúa creciendo en 2023,''
Adjust Mobile Measurement, 2023. [En línea].
Disponible en: \url{https://www.adjust.com/es/blog/finance-app-usage/}

\bibitem{OECD_INFE2023}
OECD, ``OECD/INFE 2023 International Survey of Adult Financial Literacy,''
OECD Publishing, París, 2023.
doi: 10.1787/56003a32-en

\bibitem{EncINEGI}
INEGI, ``Encuesta Nacional sobre Disponibilidad y Uso de Tecnologías de la Información en los
Hogares (ENDUTIH) 2023,'' Comunicado de Prensa 372/24, jun. 2024. [En línea].
Disponible en: \url{https://www.inegi.org.mx/contenidos/saladeprensa/boletines/2024/ENDUTIH/ENDUTIH_23.pdf}

\bibitem{Tesseract}
N. Bispo, ``Tesseract OCR Guide: Exploring Capabilities and Performance,''
Unstract Blog, 2025. [En línea].
Disponible en: \url{https://unstract.com/blog/guide-to-optical-character-recognition-with-tesseract-ocr/}

\bibitem{GoogleCloudVisionAPI}
Google LLC., ``Cloud Vision API: Derive insights from your images in the cloud,''
Google Cloud Documentation, 2024. [En línea].
Disponible en: \url{https://cloud.google.com/vision/docs/features-list}

\bibitem{Metodologias}
R. S. Pressman and B. R. Maxim, \textit{Software Engineering: A Practitioner's Approach},
9\textsuperscript{a} ed. New York, NY: McGraw-Hill, 2020.

\bibitem{ManifestoAgil}
K. Beck \textit{et al.}, ``Manifesto for Agile Software Development,'' 2001.
[En línea]. Disponible en: \url{https://agilemanifesto.org}

\bibitem{ArquitecturaCapas}
F. Buschmann and K. Henney, ``Pattern-Oriented Software Architecture,''
\textit{ResearchGate}, 2008. [En línea].
Disponible en: \url{https://www.researchgate.net/publication/228722744}

\bibitem{WhyFlutter}
I. Lotarev, ``Flutter vs Kotlin: ¿Cuál elegir para tu proyecto?,''
Adapty Blog, 2025. [En línea].
Disponible en: \url{https://adapty.io/blog/flutter-vs-kotlin/}

\bibitem{NextJSOfficial}
Vercel, ``Next.js Documentation: Introduction,''
Vercel Inc., 2025. [En línea].
Disponible en: \url{https://nextjs.org/docs}

\bibitem{WhyMySQL}
Oracle, ``MySQL: The world's most popular open source database,''
Oracle Corporation, 2024. [En línea].
Disponible en: \url{https://www.mysql.com/why-mysql/}

\bibitem{BBVA_Finanzas}
BBVA México, ``¿Qué son las finanzas personales y su importancia?,''
BBVA, nov. 2024. [En línea].
Disponible en: \url{https://www.bbva.mx/educacion-financiera/banca-digital/cuenta-digital-que-son-finanzas-personales.html}

\end{thebibliography}
\endgroup

\end{document}
